%% ====================================================================
%% TP2 MRI, FASI, UPC : comparaison de performance entre technologies
%% Compilation : pdfLaTeX, deux passes (Overleaf : compilateur par défaut)
%% Toutes les valeurs proviennent de bench/results/results_2026-11-16T16-50-50-748Z.csv
%% (agrégées par bench/aggregate.ts, fichier bench/results/summary.md).
%% ====================================================================
\documentclass[conference,a4paper]{IEEEtran}

\usepackage[T1]{fontenc}
\usepackage[utf8]{inputenc}
\usepackage[french]{babel}
\usepackage{mathptmx}
\usepackage{graphicx}
\usepackage{booktabs}
\usepackage{amsmath}
\usepackage{url}
\usepackage{pgfplots}
\pgfplotsset{compat=1.18}
\usepackage[hidelinks]{hyperref}

\renewcommand{\IEEEkeywordsname}{Mots-clés}
\renewcommand{\abstractname}{Résumé}

\begin{document}

\title{Rust et Go plus rapides que JavaScript ?\\
Quatre backends web mesurés sur une même API de tâches adossée à SQLite}

\author{
\IEEEauthorblockN{Camille MOKOLI NGOY}
\IEEEauthorblockA{FASI, UPC\\
camillemokoli@gmail.com}
\and
\IEEEauthorblockN{Joseph MBELE KALENGA}
\IEEEauthorblockA{FASI, UPC\\
tyronnklg@gmail.com}
\and
\IEEEauthorblockN{Nicolas MALEKA LIKEMBO}
\IEEEauthorblockA{FASI, UPC\\
reymaleka000@gmail.com}
\and
\IEEEauthorblockN{Béliana MALANDALA LUNGUNGU}
\IEEEauthorblockA{FASI, UPC\\
{[ADRESSE À COMPLÉTER]}}
}

\maketitle

\begin{abstract}
On lit souvent qu'un service web écrit en Rust ou en Go tient mieux la
charge qu'un service écrit en JavaScript. Nous avons voulu vérifier cette
affirmation sur un cas concret : une API REST de gestion de tâches,
adossée à SQLite, écrite quatre fois avec ElysiaJS et Hono (tous deux sur
le moteur Bun), Axum (Rust) et la bibliothèque standard net/http de Go.
Les quatre versions suivent un contrat commun et renvoient des réponses
identiques, ce que nous avons vérifié automatiquement. Chaque version a
été soumise à k6 dans un conteneur limité à 2 cœurs et 2 Gio de mémoire,
sur deux scénarios (lecture d'une tâche, mélange lecture et écriture) et
trois niveaux de charge (10, 100 et 500 utilisateurs virtuels), avec
4 essais par configuration, soit 96 mesures et 11,7 millions de requêtes.
Le résultat contredit notre hypothèse de départ. À 500 utilisateurs en
scénario mixte, Elysia traite 14\,477 requêtes par seconde et Hono
14\,000, contre 8\,953 pour Axum et 7\,371 pour Go, dont le 99\textsuperscript{e}
centile de latence atteint 308~ms. Rust reste le plus sobre en mémoire.
Nous attribuons l'essentiel de ces écarts à l'accès à la base, sérialisé
sur une seule connexion, plutôt qu'au langage lui-même. Ces résultats
valent pour cette architecture et cette machine ; ils ne permettent pas
de classer les langages en général.
\end{abstract}

\begin{IEEEkeywords}
test de charge, performance des backends web, Bun, Rust, Go, SQLite, k6
\end{IEEEkeywords}

%% ====================================================================
\section{Introduction}
%% ====================================================================
Le temps de réponse d'un service web compte pour l'utilisateur, qui
abandonne une page lente, et pour celui qui l'héberge, car un serveur
capable de traiter plus de requêtes avec le même processeur coûte moins
cher. Le choix d'une technologie backend se fait pourtant souvent sur
des réputations : Rust serait plus rapide que Go, Go plus rapide que
JavaScript, et un framework « tiendrait mieux la charge » qu'un autre.

Ces réputations s'appuient en partie sur des mesures publiques. Le projet
TechEmpower Framework Benchmarks compare depuis 2013 des centaines de
frameworks sur des tests normalisés, dont la sérialisation JSON et les
requêtes en base~\cite{techempower}. Des retours d'expérience industriels
pèsent aussi. Discord a réécrit en Rust son service « Read States »,
écrit en Go, parce que le ramasse-miettes de Go parcourait périodiquement
un très grand cache LRU en mémoire et provoquait des pics de latence ;
la version Rust s'est montrée meilleure en latence, en CPU et en
mémoire~\cite{discord}. Ces résultats portent sur des contextes précis :
une machine puissante et un code optimisé par les communautés pour
TechEmpower, un tas de plusieurs millions d'objets chez Discord. Rien
n'indique qu'ils se transposent à une petite API sur une machine modeste.

Notre question de recherche est donc la suivante : pour une même API REST
de gestion de tâches adossée à SQLite, exécutée sur 2 cœurs et 2 Gio de
mémoire, quelle technologie parmi Elysia, Hono, Axum et Go offre le
meilleur débit et les latences les plus basses, et à quel coût en CPU et
en mémoire ?

Nous avions formulé trois hypothèses avant les mesures.
H1 : Axum et Go, compilés et capables d'utiliser plusieurs cœurs,
dépassent nettement les deux versions Bun en débit.
H2 : les écritures, sérialisées par SQLite, réduisent les écarts entre
technologies dans le scénario mixte.
H3 : le tas de notre application étant minuscule, le ramasse-miettes de
Go ne produit pas l'allongement de la latence de queue observé par
Discord.

La section~II décrit l'application, les technologies, les métriques et
le protocole. La section~III explique comment les versions ont été
écrites et comment leur équivalence a été vérifiée. La section~IV
présente et discute les résultats, la section~V conclut.

%% ====================================================================
\section{Méthodologie}
%% ====================================================================
\subsection{Application de référence}
L'application est une API REST de gestion de tâches (« todo »). Une tâche
possède un identifiant entier, un titre obligatoire, une description
facultative, un indicateur d'achèvement et une date de création. Les
données sont stockées dans une table SQLite unique dont le schéma est
identique pour toutes les versions. Six routes sont exposées :
\texttt{GET /health}, \texttt{GET /tasks} (les 100 premières tâches,
triées par identifiant), \texttt{GET}, \texttt{PUT} et \texttt{DELETE}
sur \texttt{/tasks/\{id\}}, et \texttt{POST /tasks}.

Le comportement attendu est fixé par un contrat OpenAPI commun. Un
identifiant non numérique donne une erreur 400 \texttt{invalid\_id}, un
corps invalide une erreur 400 \texttt{validation\_error}, une tâche
absente une erreur 404 \texttt{not\_found}. Une création réussie renvoie
201 avec la tâche, une suppression 204 sans corps. \texttt{PUT} remplace
la tâche entière. Au démarrage, chaque version crée 1\,000 tâches
déterministes (titre « Task $i$ », achevée si $i$ est multiple de 5).

\subsection{Technologies comparées}
Nous avons retenu quatre technologies que les membres du groupe
utilisent déjà (tableau~\ref{tab:techno}). Elysia et Hono représentent
l'écosystème JavaScript et TypeScript récent ; les deux tournent sur Bun,
ce qui permet de séparer l'effet du framework de celui du moteur
d'exécution. Axum, sur le moteur asynchrone Tokio, représente Rust. Pour
Go, nous avons choisi la bibliothèque standard \texttt{net/http}, sans
framework, parce que c'est l'usage le plus courant de ce langage pour un
service HTTP.

\begin{table}[htbp]
\caption{Technologies comparées}
\label{tab:techno}
\centering
\begin{tabular}{@{}lll@{}}
\toprule
\textbf{Technologie} & \textbf{Langage} & \textbf{Version} \\
\midrule
Elysia (Bun 1.4.4) & TypeScript & 1.4.32 \\
Hono (Bun 1.4.4) & TypeScript & 4.13.18 \\
Axum (Tokio 1.54.1) & Rust & 0.8.9 \\
net/http & Go & 1.27 \\
\bottomrule
\end{tabular}
\end{table}

L'accès à SQLite passe par le pilote le plus courant de chaque
écosystème : \texttt{bun:sqlite}, intégré à Bun, pour les deux versions
TypeScript ; \texttt{rusqlite} 0.32.1, qui embarque la bibliothèque C
de SQLite, pour Rust ; \texttt{modernc.org/sqlite}, une traduction de
SQLite en Go pur qui évite CGO, pour Go. Toutes les versions utilisent
une seule connexion par processus, le journal WAL~\cite{sqlitewal} et
\texttt{synchronous = NORMAL}.

\subsection{Métriques}
Le débit, exprimé en requêtes par seconde (RPS), est le nombre de
réponses reçues divisé par la durée de la mesure. Il indique combien de
travail un serveur accomplit avec les ressources dont il dispose. k6 le
calcule à partir du compteur \texttt{http\_reqs}.

La latence est le temps entre l'envoi d'une requête et la réception
complète de la réponse, en millisecondes, mesuré par k6
(\texttt{http\_req\_duration}). Nous retenons la médiane (P50) et les
centiles 90, 95 et 99. La moyenne masque les requêtes lentes ; le P99
dit ce que subit un utilisateur sur cent, et c'est souvent lui qui
départage deux serveurs au débit voisin.

Le taux d'erreur est la part de requêtes en échec, au sens de k6 (erreur
réseau ou statut HTTP 4xx et 5xx), complétée par des vérifications du
statut attendu pour chaque requête. Un serveur rapide qui échoue n'est
pas performant.

La consommation de CPU est relevée par \texttt{docker stats} pendant la
mesure, en pourcentage d'un cœur : 100~\% correspond à un cœur occupé,
200~\% au maximum autorisé. La mémoire est le maximum de mémoire utilisée par le
conteneur pendant la mesure, en Mio. Ces deux grandeurs
disent ce que coûte le débit obtenu.

\subsection{Outils et environnement de test}
La charge est générée par k6~\cite{k6} en boucle fermée : chaque
utilisateur virtuel envoie une requête, attend la réponse et recommence
sans pause. Chaque version tourne dans un conteneur Docker limité par
\texttt{--cpus 2} et \texttt{--memory 2g}, ce qui reproduit une petite
machine virtuelle de 2 cœurs et 2 Gio. k6 s'exécute hors du conteneur,
sur la même machine hôte, un ordinateur portable sous Windows avec
Docker Desktop [MODÈLE DU PROCESSEUR ET MÉMOIRE DE L'HÔTE À COMPLÉTER].

Deux scénarios ont été mesurés. Dans \texttt{get\_one}, chaque requête
lit une tâche tirée au hasard parmi les 1\,000 semées. Le scénario
\texttt{mixed} imite un usage réaliste : 50~\% de lectures d'une tâche,
20~\% de listes, 15~\% de créations, 10~\% de modifications et 5~\% de
couples création puis suppression.

\subsection{Déroulement des tests}
Le banc d'essai est entièrement automatisé par le script
\texttt{bench/run.ts} du dépôt. Pour chaque mesure, il démarre un
conteneur neuf, donc une base fraîche de 1\,000 tâches, attend que
\texttt{/health} réponde, puis lance 3~s de chauffe (scénario mixte,
10 utilisateurs) dont les résultats sont écartés. Il lance ensuite la
mesure proprement dite pendant 8~s, échantillonne CPU et mémoire en
parallèle, enregistre une ligne dans le fichier CSV de résultats et
détruit le conteneur.

Les mesures sont organisées en 4 tours. Un tour parcourt les
2 scénarios, les 3 niveaux de charge (10, 100 et 500 utilisateurs) et,
pour chaque palier, les 4 technologies dans un ordre tiré au hasard,
afin qu'aucune ne soit systématiquement mesurée en premier ou en
dernier. Le tout représente 96 mesures. Un autre groupe peut refaire
l'expérience avec trois commandes : \texttt{bun bench/run.ts build},
\texttt{bun bench/run.ts smoke}, puis
\texttt{bun bench/run.ts run --runs 4 --scenarios get\_one,mixed
--levels 10,100,500}.

%% ====================================================================
\section{Développement}
%% ====================================================================
\subsection{Réalisation des versions}
Les quatre versions sont rassemblées dans un dépôt Git organisé en
monorepo avec Turborepo, une application par dossier sous
\texttt{apps/}. Chacune tient en un seul fichier source et suit la même
structure : ouverture de la base et réglages, création du schéma et des
données si la table est vide, préparation des requêtes SQL, puis
déclaration des routes. Les requêtes SQL sont identiques d'une version à
l'autre ; les écritures utilisent \texttt{RETURNING} pour renvoyer la
tâche sans seconde requête.

La validation des entrées est écrite à la main dans les quatre versions,
avec les mêmes règles, plutôt que confiée aux mécanismes propres à
chaque framework, qui n'auraient pas eu le même coût ni produit
exactement les mêmes erreurs. En Rust, la connexion est partagée entre
les tâches Tokio par un \texttt{Mutex} ; en Go, le pool de
\texttt{database/sql} est limité à une connexion.

Les versions, le banc d'essai et le script d'agrégation ont été écrits
avec l'aide de l'assistant Claude (Anthropic), comme la consigne le
permet ; le même assistant a servi à la rédaction de ce rapport. Le
groupe a relu le code, construit et exécuté les versions, et réalisé
toutes les mesures sur sa propre machine.

\subsection{Équivalence fonctionnelle}
Le script \texttt{bench/smoke.ts} envoie à chaque version, sur une base
fraîche, la même suite de 27 requêtes : santé, liste, lectures
présentes, absentes et mal formées, créations valides et invalides (titre
manquant, vide ou non textuel, types erronés, JSON invalide, tableau),
modifications, suppressions. Il vérifie le statut, le type de contenu et
le contenu de chaque réponse, puis enregistre une transcription
normalisée où seule la date de création est masquée. La commande
\texttt{smoke} compare enfin les transcriptions des quatre versions,
qui doivent être identiques octet pour octet. Le banc refuse de lancer
les mesures tant que ce contrôle échoue.

\subsection{Difficultés rencontrées}
Le premier contrôle d'équivalence a échoué pour Go alors que chaque
version passait ses propres vérifications. Le pilote
\texttt{modernc.org/sqlite} convertit les colonnes déclarées
\texttt{DATETIME} en valeurs temporelles Go, que la sérialisation JSON
écrit au format RFC~3339, différent du format texte renvoyé par les
autres versions. Une conversion explicite en texte dans la requête SQL a
corrigé l'écart, et le contrôle vérifie désormais le format de cette
date. Hono, d'abord prévu sur Node.js, a été exécuté sur Bun comme
Elysia : cela simplifie la comparaison entre frameworks, mais retire
Node.js de l'étude. Le démarrage de Docker Desktop sous Windows a aussi
demandé quelques réglages avant les premières mesures.

%% ====================================================================
\section{Résultats et discussion}
%% ====================================================================
\subsection{Données recueillies}
Les 96 mesures ont abouti. Elles totalisent 11\,726\,589 requêtes. Le
fichier \texttt{bench/results/results\_2026-11-16T16-50-50-748Z.csv}
contient une ligne par mesure (technologie, scénario, charge, essai,
débit, latences, erreurs, CPU, mémoire), et le dossier
\texttt{bench/results/raw/} conserve le résumé brut de k6 pour chacune.
Les valeurs présentées ici sont des moyennes sur les 4 essais ; l'écart
type du débit est indiqué. Une seule requête a échoué sur l'ensemble,
pendant un essai de Go en scénario mixte à 500 utilisateurs : le taux
d'erreur est donc nul en pratique pour toutes les versions, et nous ne
l'utilisons pas pour les départager. Le tableau~\ref{tab:resultats}
résume la charge la plus forte.

\begin{table}[htbp]
\caption{Synthèse des mesures à 500 utilisateurs (moyenne sur 4 essais)}
\label{tab:resultats}
\centering
\setlength{\tabcolsep}{3.5pt}
\begin{tabular}{@{}lrrrrr@{}}
\toprule
\textbf{Technologie} & \textbf{Débit} & \textbf{P50} & \textbf{P99} & \textbf{CPU} & \textbf{RAM} \\
 & (req/s) & (ms) & (ms) & (\%) & (Mio) \\
\midrule
\multicolumn{6}{@{}l}{\textit{Scénario get\_one}} \\
Elysia & 29\,206 $\pm$ 618 & 14,7 & 38,7 & 72 & 30 \\
Hono & 28\,190 $\pm$ 1\,122 & 15,5 & 39,9 & 64 & 29 \\
Axum & 24\,945 $\pm$ 818 & 18,3 & 40,3 & 115 & 21 \\
Go & 14\,061 $\pm$ 355 & 24,2 & 167,3 & 110 & 34 \\
\midrule
\multicolumn{6}{@{}l}{\textit{Scénario mixed}} \\
Elysia & 14\,477 $\pm$ 344 & 31,0 & 73,9 & 70 & 33 \\
Hono & 14\,000 $\pm$ 538 & 32,4 & 75,3 & 67 & 30 \\
Axum & 8\,953 $\pm$ 180 & 49,9 & 106,1 & 91 & 25 \\
Go & 7\,371 $\pm$ 187 & 46,8 & 307,5 & 105 & 37 \\
\bottomrule
\end{tabular}
\end{table}

\subsection{Analyse}
À faible charge, les quatre versions se ressemblent. À 10 utilisateurs
en lecture, Elysia, Hono et Axum se situent entre 11\,882 et
12\,082 requêtes par seconde, des écarts inférieurs à l'écart type de
leurs propres essais, et leurs médianes sont toutes proches de 0,6~ms.
Go est un peu en retrait, à 10\,547 requêtes par seconde.

La figure~\ref{fig:barres} montre que les écarts apparaissent dès
100 utilisateurs. En lecture, Elysia atteint 29\,789 requêtes par
seconde et Hono 28\,586, Axum 23\,714, et Go 14\,308, soit environ deux
fois moins que les versions Bun.

\begin{figure}[htbp]
\centering
\begin{tikzpicture}
\begin{axis}[
  width=\columnwidth, height=5cm,
  ybar, bar width=14pt, ymin=0, ymax=35000,
  symbolic x coords={Elysia,Hono,Axum,Go},
  xtick=data,
  ylabel={Débit (requêtes/s)},
  nodes near coords,
  nodes near coords style={font=\scriptsize},
  scaled y ticks=false,
  yticklabel style={/pgf/number format/1000 sep={\,}},
  tick label style={font=\footnotesize},
  label style={font=\footnotesize},
  enlarge x limits=0.18]
\addplot[fill=gray!40, draw=black] coordinates
  {(Elysia,29789) (Hono,28586) (Axum,23714) (Go,14308)};
\end{axis}
\end{tikzpicture}
\caption{Débit moyen sur 4 essais, scénario get\_one, 100 utilisateurs simultanés.}
\label{fig:barres}
\end{figure}

La figure~\ref{fig:debit} suit le débit du scénario mixte selon la
charge. Les versions Bun passent d'environ 10\,000 à 14\,500 requêtes
par seconde entre 10 et 100 utilisateurs, puis plafonnent. Axum et Go
plafonnent dès 10 utilisateurs, autour de 9\,000 et 7\,400 requêtes par
seconde. Ajouter des utilisateurs au-delà de 100 n'augmente le débit
d'aucune version ; cela allonge seulement l'attente.

\begin{figure}[htbp]
\centering
\begin{tikzpicture}
\begin{semilogxaxis}[
  width=\columnwidth, height=5cm,
  xlabel={Utilisateurs simultanés (échelle logarithmique)},
  ylabel={Débit (requêtes/s)},
  xtick={10,100,500}, xticklabels={10,100,500},
  scaled y ticks=false,
  yticklabel style={/pgf/number format/1000 sep={\,}},
  tick label style={font=\footnotesize},
  label style={font=\footnotesize},
  legend style={font=\scriptsize, at={(0.03,0.97)}, anchor=north west},
  ymin=0, ymax=17000]
\addplot[black, thick, mark=*] coordinates {(10,10049) (100,14644) (500,14477)};
\addplot[black, thick, dashed, mark=square*] coordinates {(10,9982) (100,14122) (500,14000)};
\addplot[black, thick, dotted, mark=triangle*] coordinates {(10,8729) (100,9324) (500,8953)};
\addplot[gray, thick, dashdotted, mark=diamond*] coordinates {(10,7654) (100,7299) (500,7371)};
\legend{Elysia, Hono, Axum, Go}
\end{semilogxaxis}
\end{tikzpicture}
\caption{Évolution du débit selon la charge, scénario mixed (moyenne sur 4 essais).}
\label{fig:debit}
\end{figure}

La latence de queue sépare plus nettement les versions
(figure~\ref{fig:p99}). À 500 utilisateurs en scénario mixte, le P99
vaut 74~ms pour Elysia, 75~ms pour Hono, 106~ms pour Axum et 308~ms pour
Go. En lecture, Axum rejoint les versions Bun (40~ms) alors que Go monte
à 167~ms, quatre fois plus. Chez Go, l'écart entre la médiane et le
P99 est le plus grand : un utilisateur sur cent attend plus de six fois
plus que l'utilisateur médian.

\begin{figure}[htbp]
\centering
\begin{tikzpicture}
\begin{loglogaxis}[
  width=\columnwidth, height=5cm,
  xlabel={Utilisateurs simultanés (échelle logarithmique)},
  ylabel={Latence P99 (ms, échelle log.)},
  xtick={10,100,500}, xticklabels={10,100,500},
  tick label style={font=\footnotesize},
  label style={font=\footnotesize},
  legend style={font=\scriptsize, at={(0.03,0.97)}, anchor=north west}]
\addplot[black, thick, mark=*] coordinates {(10,3.32) (100,20.96) (500,73.85)};
\addplot[black, thick, dashed, mark=square*] coordinates {(10,3.23) (100,21.54) (500,75.29)};
\addplot[black, thick, dotted, mark=triangle*] coordinates {(10,3.35) (100,23.12) (500,106.10)};
\addplot[gray, thick, dashdotted, mark=diamond*] coordinates {(10,3.77) (100,64.79) (500,307.50)};
\legend{Elysia, Hono, Axum, Go}
\end{loglogaxis}
\end{tikzpicture}
\caption{Latence au 99\textsuperscript{e} centile selon la charge, scénario mixed (moyenne sur 4 essais).}
\label{fig:p99}
\end{figure}

Le coût en ressources inverse en partie ce classement. Les versions Bun
n'ont jamais dépassé 92~\% d'un cœur, alors qu'Axum et Go ont occupé
jusqu'à plus de 1,5 cœur. Pour le même travail, Axum utilise le moins de
mémoire : 6 à 25~Mio selon la charge, contre 14 à 37~Mio pour Go et
20 à 33~Mio pour les versions Bun. Enfin, Elysia et Hono ne se
distinguent pas : à chaque palier, leur écart est du même ordre que la
dispersion de leurs essais.

\subsection{Interprétation et discussion}
L'hypothèse H1 est réfutée dans nos conditions. Les versions Bun sont
les plus rapides alors qu'elles exécutent le code de l'application sur
un seul fil, comme le montre leur CPU plafonné sous un cœur. Notre
explication tient à l'accès à la base. Toutes les versions partagent une
connexion SQLite unique. Sous Bun, les requêtes passent l'une après
l'autre sur le même fil, sans verrou ni changement de contexte. Axum et
Go répartissent les requêtes sur plusieurs fils, qui se disputent ensuite
cette unique connexion : le verrou de Rust et la file d'attente du pool
de Go ajoutent un coût de synchronisation sans apporter de parallélisme,
puisque le travail utile reste sérialisé. Le CPU supplémentaire consommé
par Axum et Go, sans gain de débit, va dans le même sens.

Le retard propre à Go s'explique probablement aussi par son pilote :
\texttt{modernc.org/sqlite} est une traduction automatique de SQLite en
Go, alors que \texttt{bun:sqlite} et \texttt{rusqlite} appellent la
bibliothèque C d'origine. Notre étude compare donc des couples langage
et pilote, et non des langages isolés. Nous n'avons pas fait de
profilage, et ces explications restent des hypothèses à vérifier.

L'hypothèse H2 n'est confirmée qu'en partie. Le scénario mixte divise
bien le débit de toutes les versions, mais il ne rapproche pas les
technologies : le rapport entre Elysia et Go reste voisin de deux, et
Axum perd davantage que les autres (il conserve 36~\% de son débit de
lecture, contre environ 50~\% pour les autres versions). La liste de
100 tâches et les écritures, qui gardent le verrou plus longtemps,
pourraient en être la cause.

Pour H3, nos données ne permettent pas de conclure. Go présente bien la
latence de queue la plus élevée, mais nous n'avons relevé que des
centiles agrégés sur 8~s, sans trace du ramasse-miettes. L'attente dans
la file de la connexion unique suffit à expliquer un P99 qui croît avec
le nombre d'utilisateurs, sans qu'il soit nécessaire d'invoquer le
ramasse-miettes. La situation de Discord, un tas de plusieurs millions
d'objets~\cite{discord}, n'a d'ailleurs rien de commun avec la nôtre.

Plusieurs limites encadrent ces résultats. Quatre essais par
configuration, au lieu des quarante de l'exemple du cours, donnent des
moyennes stables (écarts types inférieurs à 5~\% du débit pour la
plupart) mais une estimation fragile de la dispersion. k6 et Docker
partagent la même machine : à forte charge, le générateur de charge
consomme une part du processeur et peut plafonner lui-même, ce qui
pourrait limiter les versions les plus rapides. Les latences maximales
proches d'une seconde, observées pour toutes les versions en lecture à
500 utilisateurs, ressemblent à un délai de retransmission réseau du
système hôte plutôt qu'à un comportement des serveurs. Nous n'avons pas
mesuré de scénario purement écriture, ni d'autres réglages de connexion.
Enfin, chaque mesure ne dure que 8~s ; des effets plus lents, comme la
croissance de la mémoire, échappent à ce protocole.

%% ====================================================================
\section{Conclusion}
%% ====================================================================
Nous nous demandions quelle technologie, parmi Elysia, Hono, Axum et Go,
sert le mieux une API de tâches adossée à SQLite sur une machine de
2 cœurs et 2 Gio. Dans nos mesures, les deux versions Bun offrent le
meilleur débit et les latences de queue les plus basses, Axum vient
ensuite avec la plus faible empreinte mémoire, et Go ferme la marche.
Ce classement va à l'encontre de la réputation des langages compilés,
et nous l'attribuons d'abord à la conception de l'accès à la base, une
connexion unique qui annule l'avantage du multi-cœur, ainsi qu'aux
différences entre pilotes SQLite.

Le résultat le plus utile de cette étude est peut-être celui-là : sur
une application de ce type, l'accès aux données pèse davantage que le
langage du serveur. Pour aller plus loin, il faudrait répéter les
mesures sur une machine dédiée au générateur de charge, tester un pool
de connexions en lecture pour Axum et Go, remplacer
\texttt{modernc.org/sqlite} par un pilote C en Go, et profiler chaque
version pour confirmer les explications avancées ici.

%% ====================================================================
\begin{thebibliography}{00}
\bibitem{techempower} TechEmpower, « Framework Benchmarks Round 23 »,
2025. [En ligne]. Disponible :
\url{https://www.techempower.com/blog/tag/tfb/}. [Consulté le 16 novembre 2026].
\bibitem{discord} J. Howarth, « Why Discord is switching from Go to
Rust », Discord Blog, février 2020. [En ligne]. Disponible :
\url{https://discord.com/blog/why-discord-is-switching-from-go-to-rust}.
[Consulté le 16 novembre 2026].
\bibitem{sqlitewal} SQLite, « Write-Ahead Logging ». [En ligne].
Disponible : \url{https://www.sqlite.org/wal.html}. [Consulté le
16 novembre 2026].
\bibitem{k6} Grafana Labs, « Grafana k6 documentation ». [En ligne].
Disponible : \url{https://grafana.com/docs/k6/latest/}. [Consulté le
16 novembre 2026].
\bibitem{tp} Faculté des sciences informatiques, UPC, « Méthode de
recherche en informatique, travail pratique n\textsuperscript{o}~2 :
comparer la performance d'une même application écrite avec plusieurs
technologies », novembre 2026.
\end{thebibliography}

\end{document}
